%% kapa: exam=1 round=19 year=2008 period=1 form=A subject=accounting subject_ko=회계학 seq=16 number=96
%% source: Q-Net 감정평가사 기출문제 (cst00309) artlSeq=1001962 file=extracted/1008/1001962/2008년 제19회 감정평가사 1차시험 기출문제/1교시.hwp
%% conversion: kordoc 4.24.51 -> tools/gen1.py (auto); review=proofread-sonnet; flags=none
\begin{question}{96}
(주)단결은 20X7년 1월 1일에 액면이자율 연 9\%, 5년 만기, 액면금액 ￦600,000의 사채를 ￦576,834에 발행하였다. 이 발행가액은 10\%의 시장이자율을 반영하고 있다. 이 사채의 이자지급일은 1월 1일과 7월 1일로 연 2회이다. (주)단결이 유효이자율법을 적용하여 사채할인발행차금을 상각한다면, 20X7년 7월 1일의 이자지급 시 상각되는 사채할인발행차금은 얼마인가? (단, 사채이자는 월할로 계산하고, 소수점 이하 첫째자리에서 반올림하시오.)
\begin{choices}
  \item ￦1,842
  \item ￦2,085
  \item ￦2,317
  \item ￦3,683
  \item ￦23,166
\end{choices}
\end{question}
\answer{1}
